\begin{lemma}[Action of the energy-normalized gradient on Hermite functions]
Let $D = 2^{-1/2}\nabla$, so that $D_j = 2^{-1/2}\partial_j$.  For every
multi-index $\alpha \in \mathbb{N}_0^n$, every $j \in \{1,\ldots,n\}$ and every
$x \in \mathbb{R}^n$, the Hermite function $h_\alpha$ of \noderef{HermiteFun} is
differentiable at $x$ and
\[
  \partial_j h_\alpha(x) = \sqrt{2}\,\sqrt{\alpha_j}\; h_{\alpha - e_j}(x),
  \qquad\text{equivalently}\qquad
  D_j h_\alpha = \sqrt{\alpha_j}\, h_{\alpha - e_j},
\]
where $\alpha - e_j$ is the multi-index of \noderef{MultiIndexSum} obtained by
decreasing the $j$th entry of $\alpha$ by one, with the convention that this entry
stays at $0$ when $\alpha_j = 0$ (in which case the right-hand side vanishes
anyway, because of the factor $\sqrt{\alpha_j} = 0$).
\end{lemma}
\begin{proof}
\emph{Step 1: the one-dimensional identity.}
We first prove
\begin{equation*}
  h_k'(s) = \sqrt2\,\sqrt{k}\; h_{k-1}(s),
  \qquad k \in \mathbb{N}_0,\ s \in \mathbb{R}, \tag{1}
\end{equation*}
with the right-hand side read as $0$ when $k = 0$.

For $k = 0$: by \noderef{Hermite1Zero}, $h_0 \equiv 1$, so $h_0' = 0$; and the
right-hand side carries the factor $\sqrt{0} = 0$.  Both sides vanish.

For $k \geq 1$ write $k = m+1$.  By \noderef{Hermite1},
$h_{m+1}(s) = \mathrm{He}_{m+1}(\sqrt2\,s)/\sqrt{(m+1)!}$, a polynomial in $s$, hence
differentiable, and the chain rule gives
\[
  h_{m+1}'(s) = \frac{\sqrt2\;\mathrm{He}_{m+1}'(\sqrt2\,s)}{\sqrt{(m+1)!}} .
\]
By \noderef{HermiteDerivative}, $\mathrm{He}_{m+1}' = (m+1)\,\mathrm{He}_m$, so
\[
  h_{m+1}'(s)
    = \frac{\sqrt2\,(m+1)\,\mathrm{He}_m(\sqrt2\,s)}{\sqrt{(m+1)!}}
    = \sqrt2\,\frac{m+1}{\sqrt{m+1}}\cdot
        \frac{\mathrm{He}_m(\sqrt2\,s)}{\sqrt{m!}}
    = \sqrt2\,\sqrt{m+1}\; h_m(s),
\]
where the middle step used $(m+1)! = (m+1)\,m!$, so that
$\sqrt{(m+1)!} = \sqrt{m+1}\,\sqrt{m!}$, and the last step used
\noderef{Hermite1} again.  This is (1) with $k = m+1$.

\emph{Step 2: differentiability.}
By \noderef{HermiteFunContDiff}, $h_\alpha \in C^\infty(\mathbb{R}^n)$.  A
$C^\infty$ function is in particular differentiable at every point, so
$h_\alpha$ is differentiable at $x$; this is the first assertion, and it also
licenses the use of $\partial_j h_\alpha(x)$ and of the product rule below.

\emph{Step 3: the $n$-dimensional identity.}
By \noderef{HermiteFun}, $h_\alpha(x) = \prod_{i=1}^n h_{\alpha_i}(x_i)$.  Only
the $i = j$ factor depends on $x_j$, so the product rule gives
\[
  \partial_j h_\alpha(x)
    = h_{\alpha_j}'(x_j) \prod_{i \neq j} h_{\alpha_i}(x_i) .
\]
Substituting (1) with $k = \alpha_j$ and $s = x_j$,
\[
  \partial_j h_\alpha(x)
    = \sqrt2\,\sqrt{\alpha_j}\;
      h_{\alpha_j - 1}(x_j) \prod_{i \neq j} h_{\alpha_i}(x_i) .
\]
The multi-index $\alpha - e_j$ agrees with $\alpha$ in every entry $i \neq j$ and
has $j$th entry $\alpha_j - 1$, so by \noderef{HermiteFun} the product on the
right is exactly $h_{\alpha-e_j}(x)$.  Hence
$\partial_j h_\alpha(x) = \sqrt2\sqrt{\alpha_j}\, h_{\alpha-e_j}(x)$.  When
$\alpha_j = 0$ the prefactor is $0$, so the identity holds regardless of how
$\alpha - e_j$ is interpreted, which justifies the stated convention.  Dividing
by $\sqrt2$ gives the form for $D_j = 2^{-1/2}\partial_j$.
\end{proof}
